%Pretex 2/16/99 lmd
% CAH 1st pass 9/7/99
% ADG Finals 1.20.99

\documentclass[final]{siamltex}
%\usepackage{amsfonts}
\usepackage{amssymb}

%\advance\oddsidemargin by -1.1in
%\advance\evensidemargin by -1.1in
%\advance\topmargin by -.9in
%\advance\textheight by 1.5in
%\def\baselinestretch{1.5}
%\advance\textwidth by 2.2in
%\pagestyle{headings}
%\newtheorem{theorem}{Theorem}
%\newtheorem{corollary}{Corollary}
%\newtheorem{lemma}{Lemma}
\newcommand{\arw}{\mbox{$\rightarrow$}}
\newcommand{\vs}{\vspace{.15in}}
\newcommand{\vb}{\vspace{.25in}}
\newcommand{\vbb}{\vspace{.4in}}


\title{An Inequality Involving the Generalized Hypergeometric Function
and the Arc Length
\newline
of an Ellipse\thanks{Received by the editors July 7, 1998; 
accepted for publication (in revised form) December 17, 1998; published electronically
March 6, 2000. \URL sima/31-3/34157.html}}

\author{Roger W. Barnard\thanks{Department of 
Mathematics and Statistics,  Texas Tech University, 
Lubbock, TX 79409 
(barnard@math.ttu.edu, pearce@math.ttu.edu).}
\and Kent Pearce\footnotemark[2] 
\and Kendall C. Richards\thanks {Department of 
Mathematics, Southwestern University, 
Georgetown, TX  78626 (richards@\allowbreak southwestern.edu).}}

\begin{document}

\maketitle
\vspace{-1.2in}
\slugger{sima}{2000}{31}{3}{693--699}
\vspace{.9in}

\setcounter{page}{693}

\begin{abstract}
In this paper we verify a conjecture of M. Vuorinen that the
Muir approximation is a lower approximation to the arc length of
an ellipse. Vuorinen conjectured that 
$f(x)={}_{2}F_{1}({\frac{1}{2}},-{\frac{1}{2}};1;x)-[(1+(1-x)^{3/4})/2]^{2/3}$ is positive for 
$x\in (0,1)$. The authors prove a much stronger result which says that the Maclaurin coefficients of $f$
are nonnegative. As a key lemma, we show that 
${}_{3}F_{2}(-n,a,b;1+a+b,1+\epsilon -n;1)>0$ when $0<ab/(1+a+b)<\epsilon <1$ for all positive integers 
$n$.
\end{abstract}

\begin{keywords} 
hypergeometric, approximations, elliptical arc length 
\end{keywords}

\begin{AMS}
33C, 41A
\end{AMS}

\begin{PII}
S0036141098341575
\end{PII}

\pagestyle{myheadings}
\thispagestyle{plain}
%\markboth{R. W. BARNARD, K. PEARCE, AND K. C. RICHARDS}{GENERALIZED HYPERGEOMETRIC FUNCTIONS AND ELLIPTICAL ARC LENGTH}
\markboth{R. W. BARNARD, K. PEARCE, AND K. C. RICHARDS}{${}_{p}F_{q}$ AND ELLIPTICAL ARC LENGTH}

\section{Introduction}
%\noindent
Let $a$ and $b$ be the semiaxes of an ellipse with eccentricity $e = \sqrt{a^2 - b^2}/a$.  Let
$L(a,b)$ denote the arc length of the ellipse.  Without loss of generality we can take one of the semiaxes, 
say $a$, to be $1$.  Legendre's complete elliptic integral of the second kind can be defined by 

\[
E(r) = \int_{0}^{\pi/2} \sqrt{ 1 - r^2 \sin^2 t}\ dt.
\]
Elliptic integrals are so named because of their connection with $L(a,b)$.  In turn, these are related to 
Gauss's hypergeometric functions, ${}_{2}F_{1}$, defined by

\[
{}_{2}F_{1}(a_1,a_2;b_1;z) = \sum _{n=0}^{\infty }\frac{ (a_1)_n (a_2)_n}{(b_1)_n n!} z^n
\]
with the Appell (or Pochhammer) symbol $(a)_{n}=a(a+1)\cdot \cdot \cdot (a+n-1)$ for $n\geq 1$ and $(a)_{0}=1, a \neq 0$.
We shall need the generalized hypergeometric function, ${}_{p}F_{q}$, defined by

\[
\,{\ }_{p}F_{q}(a_{1},a_{2},\ldots,a_{p};b_{1},b_{2},\ldots, b_{q};z)=\sum_{n=0}^{%
\infty }\frac{(a_{1})_{n}\cdot (a_{2})_{n}\cdot \cdot \cdot (a_{p})_{n}}{%
(b_{1})_{n}\cdot (b_{2})_{n}\cdot \cdot \cdot (b_{q})_{n}}\frac{z^{n}}{n!} 
\]
(see \cite[p. 73]{R}).  It was noted by Maclaurin in 1742 (see \cite{AB}) that
$$L(1,b) = 4 E(e) = 2\pi {}_{2}F_{1}(\mbox{$\frac{1}{2}$},-\mbox{$\frac{1}{2}$};1;e^2).$$
There are various references, books, and articles, which discuss the relationships between elliptic integrals and 
hypergeometric functions (see \cite{AVV}, \cite{C}) and their role in applications to physics (see \cite{L}, \cite{G}) 
and in geometric function theory  (see \cite{H}, \cite{AVV}).  From antiquity several more easily computable 
approximations to $L(a,b)$ have been suggested.  The Almkvist--Berndt survey article \cite{AB} has an extensive 
discussion of these approximations.  These approximations and their historical and recent connections to the 
approximations of $\pi$ can be found in the Borweins' book \cite{BB}.  An excellent source for all of the above 
ideas is the Anderson--Vamanamurthy--Vuorinen book {\em Conformal Invariants, Inequalities, and Quasiconformal Mappings} \cite{AVV}.

In 1883, it was proposed by Muir  (see \cite{AB}) that $L(1,b)$ could be simply approximated by 
$2\pi [(1+b^{3/2})/2]^{2/3}$. A close numerical examination of the error in this
approximation lead M. Vuorinen to pose Problem 5.6 in  \cite{V}. This was announced at
several international conferences. Letting $x=1-b^{2}$, he asked whether the
Muir approximation 
\[
g(x)=\left( \frac{1+(1-x)^{3/4}}{2}\right) ^{2/3} 
\]
is a lower approximation for the value given by the hypergeometric function 
\[
h(x)={}_{2}F_{1}\left( \mbox{$\frac{1}{2}$},-\mbox{$\frac{1}{2}$};1;x\right)
, 
\]
that is, whether 
\[
h(x)-g(x)\geq 0\;\mbox{for all}\;x\in (0,1).  
\]

We shall prove the following much stronger result.

\begin{theorem}\label{strong}
Let $g(x)=\sum_{n=0}^{\infty
}a_{n}x^{n}$ and $h(x)=\sum_{n=0}^{\infty }A_{n}x^{n}$. Then, 
\begin{equation}
a_{k}\leq A_{k}\;\mbox{for all}\;k=0,1,2,\ldots,n,\ldots.  \label{eq2}
\end{equation}
In particular, the function  $f(x) \equiv [h(x)-g(x)]/x^{4}$ is convex and increasing from 
$(0,1]$ onto $(\alpha,\beta]$, where $\alpha = 2^{-14} = 0.000061\cdots$ and 
$\beta = ({2}/{\pi })-{2^{-2/3}} = 0.006659\cdots$.
\end{theorem}

{\em Remarks.} The ideas and techniques used to prove Lemma \ref{required} and Theorem \ref{strong} will be used in 
\cite{BPR} to determine surprising hierarchical relationships among the 13 historical approximations to $L(a,b)$ 
discussed in \cite{AB}.  These approximations range over four centuries from Kepler's in 1642 to Almkvist's in 1985 
and include two from Ramanujan.   
%Deleted per AA %%% In addition, the result in Lemma \ref{required} has been extended to a wider range of 
%parameters.

\section{Proof of main results}
The proof of Theorem \ref{strong} requires the following lemma.  

%\smallskip

\begin{lemma}\label{required}
Suppose $a,b>0$. Then, for any $\epsilon $
satisfying $\frac{ab}{1+a+b}<\epsilon <1$, 
\[
\,_{3}F_{2}(-n,a,b;1+a+b,1+\epsilon -n;1)>0\ \ 
\mbox{ for all integers $n
\geq 1$.} 
\]
\end{lemma}

For the reader's convenience, we include the following classical
identities.

{\bf Identity 1} (see \cite [p. 558, eq. (15.2.24)]{AS}). If 
$|z|<1$, then 
\[
(c-b-1)\cdot {}_{2}F_{1}(a,b;c;z)=(c-1)\cdot {}_{2}F_{1}(a,b;c-1;z)-b\cdot {\ }_{2}F_{1}(a,b+1;c;z). 
\]

{\bf Identity 2} (see \cite[p. 60, Thm. 21]{R}). If $|z|<1$, then 
\[
{\ }_{2}F_{1}(a,b;c;z)=(1-z)^{c-a-b}\cdot {}_{2}F_{1}(c-a,c-b;c;z). 
\]

{\bf Identity 3} (see \cite[p. 59, eq. (3.1.1)]{GR}). If $F\,=\,{\ }_{3}F_{2}$, then 
\[
F(-n,a,b;c,d;1)=\frac{(d-b)_{n}}{(d)_{n}}F(-n,c-a,b;c,1+b-d-n;1). 
\]

{\bf Identity 4} (see \cite [p. 82, eq. (14)]{R}). If $F\,=\,{\ }_{3}F_{2}$ and $|z|<1$, then
\begin{eqnarray*}
(a_{1}-a_{2}) &\cdot & F(a_{1},a_{2},a_{3};b_{1},b_{2};z)\\&=&a_{1}\cdot
F(a_{1}+1,a_{2},a_{3};b_{1},b_{2};z)-a_{2}\cdot
F(a_{1},a_{2}+1,a_{3};b_{1},b_{2};z). 
\end{eqnarray*}

{\it Proof of Lemma~$\ref{required}$}.
Using an idea suggested in \cite
{2}, we let $F=\,_{3}F_{2}$ and consider the generating function 
\[
f(r)=\sum_{n=0}^{\infty }\frac{-(-\epsilon )_{n}}{n!}F(-n,a,b;1+a+b,1+\epsilon -n;1)r^{n}=
\sum_{n=0}^{\infty }c_{n}r^{n}, 
\]
where $|r|<1$. Note that $-(-\epsilon )_{n}>0$ for $0 < \epsilon < 1$ and for all $n\geq 1$. Thus we
seek to verify that $c_{n}>0$ for all $n\geq 1$.

%\noindent
In this direction, we have 
\begin{eqnarray*}
f(r) &=&\sum_{n=0}^{\infty }\frac{-(-\epsilon )_{n}}{n!}\sum_{k=0}^{n}\frac{%
(-n)_{k}(a)_{k}(b)_{k}}{(a+b+1)_{k}(1+\epsilon -n)_{k}k!}r^{n}  \nonumber \\
&=&\sum_{n=0}^{\infty }\frac{-(-\epsilon )_{n}}{(1)_{n}}\sum_{k=0}^{n}\frac{%
\frac{(-1)^{k}(1)_{n}}{(1)_{n-k}}(a)_{k}(b)_{k}}{(a+b+1)_{k}\frac{%
(-1)^{k}(-\epsilon )_{n}}{(-\epsilon )_{n-k}}k!}r^{n} \\ 
&\ &\ \left\{\mbox{using $(\alpha)_{n-k} = \frac{(-1)^{k}(\alpha)_{n}}{(1-\alpha-n)_{k}}$ and
$(1)_{n} = n!$}\right\} \nonumber \\
&=&-\sum_{n=0}^{\infty }\sum_{k=0}^{n}\left( \frac{(a)_{k}(b)_{k}}{%
(a+b+1)_{k}k!}r^{k}\right) \left( \frac{(-\epsilon )_{n-k}}{(n-k)!}%
r^{n-k}\right) \\
&=&-\sum_{n=0}^{\infty }\left( \frac{(-\epsilon )_{n}}{(n)!}r^{n}\right)
\sum_{k=0}^{\infty }\left( \frac{(a)_{k}(b)_{k}}{(a+b+1)_{k}k!}r^{k}\right) \ \ \ \ 
\mbox{(see \cite[p. 57, eq. (2)]{R})} \nonumber \\
&=&-(1-r)^{\epsilon }\,{\ }_{2}F_{1}(a,b;a+b+1;r).
\end{eqnarray*}
Differentiating, we have 
\begin{eqnarray}
f^{\prime }(r)&=&\epsilon (1-r)^{\epsilon -1}\,{\ }_{2}F_{1}(a,b;a+b+1;r)  \label{eqA} \\
&-&\frac{ab(1-r)^{\epsilon }}{(a+b+1)}\,{\ }_{2}F_{1}(a+1,b+1;a+b+2;r).  \nonumber
\end{eqnarray}
An application of Identity 1 followed by Identity 2 
%to $\,{\ }_{2}F_{1}(a+1,b+1;a+b+2;r)$ yields
to \hspace{-.25pc}${\ }_{2}F_{1}(a+1,b+1;a+b+2;r)$ yields
\begin{eqnarray*}
&& \frac{ab(1-r)^{\epsilon }}{(a+b+1)} {}_{2}F_{1}(a+1,b+1;a+b+2;r) \\
&&\ \ \ \ \ \ =\frac{b(1-r)^{\epsilon }}{(a+b+1)}[ (a+b+1)\cdot
{}_{2}F_{1}(a+1,b+1;a+b+1;r) \nonumber \\
&&\ \ \ \ \ \ -\; (b+1)\cdot {}_{2}F_{1}(a+1,b+2;a+b+2;r)] \nonumber \\
&&\ \ \ \ \ \ =(1-r)^{\epsilon -1}\left[ b\cdot {}_{2}F_{1}(a,b;a+b+1;r) 
- \frac{b(b+1)}{(a+b+1)}\cdot {}_{2}F_{1}(a,b+1;a+b+2;r)\right] . \nonumber
\end{eqnarray*}
Thus (\ref{eqA}) becomes 

%\newpage
\begin{eqnarray}
f^{\prime }(r) &=&(1-r)^{\epsilon -1}\Bigg[ (\epsilon -b)\cdot
{}_{2}F_{1}(a,b;a+b+1;r) \nonumber \\
%&+& \frac{b(b+1)}{(a+b+1)}\cdot
%{}_{2}F_{1}(a,b;a+b+1;r) \Bigg] \nonumber \\ 
&+& \frac{b(b+1)}{(a+b+1)}\cdot
{}_{2}F_{1}(a,b+1;a+b+2;r)\Bigg] \nonumber \\
&=&(1-r)^{\epsilon -1}\sum_{n=0}^{\infty }\frac{(a)_{n}}{n!}\left[ \frac{%
(b)_{n}(\epsilon -b)}{(a+b+1)_{n}}
+ \frac{b(b+1)(b+1)_{n}}{(a+b+1)(a+b+2)_{n}}\right] r^{n} \nonumber \\
&=&(1-r)^{\epsilon -1}\sum_{n=0}^{\infty }\frac{(a)_{n}}{n!}\left[ \frac{%
(b)_{n}(\epsilon -b)}{(a+b+1)_{n}}  
+ \frac{(b+1)(b)_{n}(b+n)}{(a+b+1)_{n}(a+b+1+n)}\right] r^{n}  \label{eqB} \\
&=&(1-r)^{\epsilon -1}\sum_{n=0}^{\infty }\frac{(a)_{n}(b)_{n}}{%
(a+b+1)_{n}(a+b+1+n)n!} \nonumber \\
&\times & \left[ (\epsilon -b)(a+b+1+n)+(b+1)(b+n)\right] r^{n} \nonumber \\
&=&(1-r)^{\epsilon -1}\sum_{n=0}^{\infty }\frac{(a)_{n}(b)_{n}}{%
(a+b+1)_{n}(a+b+1+n)n!} \label{eqC}\\ 
&\times &\left[ \epsilon (a+b+1+n)+n-ab\right] r^{n}, \nonumber
\end{eqnarray}
where (\ref{eqB}) makes use of $\alpha (\alpha +1)_{n}=(\alpha )_{n}(\alpha
+n)$. \noindent If $\frac{ab}{a+b+1}<\epsilon <1$, then the expression in (%
\ref{eqC}) is the product of two series with all positive Maclaurin series
coefficients. Hence $f^{\prime }$ has all positive Maclaurin series
coefficients which is equivalent to the desired result.
%\thinspace \rule{2.2mm}{3.3mm}
\qquad\endproof

\begin{corollary}\label{integers}
Let $T_{n}=\;{}_{3}F_{2}\left( -n,%
\mbox{$\frac{3}{2}$},\mbox{$\frac{1}{2}$};2,\mbox{$\frac{5}{4}$}-n;1\right) $%
. Then, for all integers $n\geq 8$, 
\[
T_{n+1}>T_{n}>0. 
\]
\end{corollary}

%\noindent {\bf Proof of Corollary 2.2:}\ \ 
\begin{proof}
Let $F=\,_{3}F_{2}$ and $\displaystyle %
B_{n}=\left( \mbox{$\frac{3}{4}$}-n\right) _{n}/\left( \mbox{$\frac{5}{4}$}%
-n\right) _{n}$. Using Identity 3, we have that 
\[
T_{n}=B_{n}F\left( -n,\mbox{$\frac{1}{2}$},\mbox{$\frac{1}{2}$};2,%
\mbox{$\frac{1}{4}$};1\right) . 
\]
Direct calculation reveals that $%
T_{9}>T_{8}>0>T_{7}>\cdots>T_{2}=T_{1}$. Now suppose that $T_{n}>T_{n-1}>0$ for
some $n\geq 9$ and note that $\displaystyle   B_{n+1}/B_{n}=\left( n+%
\mbox{$\frac{1}{4}$}\right) /\left( n-\mbox{$\frac{1}{4}$}\right) $. Then, 
\begin{eqnarray}
T_{n+1} &=&B_{n+1}F\left( -n-1,\mbox{$\frac{1}{2}$},%
\mbox{$\frac{1}{2}$};2,\mbox{$\frac{1}{4}$};1\right) \nonumber \\
&=&\frac{B_{n+1}}{\left( n+\mbox{$\frac{3}{2}$}\right) }\left[ (n+1)F\left(
-n,\mbox{$\frac{1}{2}$},\mbox{$\frac{1}{2}$};2,\mbox{$\frac{1}{4}$};1\right)  
+ \mbox{$\frac{1}{2}$}F\left( -n-1,\mbox{$\frac{3}{2}$},\mbox{$\frac{1}{2}$}%
;2,\mbox{$\frac{1}{4}$};1\right) \right]  \label{eqD} \\
&=&\frac{B_{n+1}(n+1)}{B_{n}\left( n+\mbox{$\frac{3}{2}$}\right) }T_{n} 
+\frac{B_{n+1}}{2\left( n+\mbox{$\frac{3}{2}$}\right) }F\left( -n-1,%
\mbox{$\frac{3}{2}$},\mbox{$\frac{1}{2}$};2,\mbox{$\frac{1}{4}$};1\right) \nonumber \\
&=&\frac{\left( n+\mbox{$\frac{1}{4}$}\right) (n+1)}{\left( n-%
\mbox{$\frac{1}{4}$}\right) \left( n+\mbox{$\frac{3}{2}$}\right) }T_{n} 
+\frac{B_{n+1}}{2\left( n+\mbox{$\frac{3}{2}$}\right) }F\left( -n-1,%
\mbox{$\frac{3}{2}$},\mbox{$\frac{1}{2}$};2,\mbox{$\frac{1}{4}$};1\right) \nonumber \\
&>&T_{n}+\frac{B_{n+1}}{2\left( n+\mbox{$\frac{3}{2}$}\right) }F\left( -n-1,% \nonumber
\mbox{$\frac{3}{2}$},\mbox{$\frac{1}{2}$};2,\mbox{$\frac{1}{4}$};1\right), \nonumber 
%&\ &\ \left\{\mbox{since
%$\frac{(n+1/4)(n+1)}{(n-1/4)(n+3/2)} >1$ and $T_{n} > 0$}\right\}, \nonumber
\end{eqnarray}
where (\ref{eqD}) follows from Identity 4, and the inequality holds because
$\frac{(n+1/4)(n+1)}{(n-1/4)(n+3/2)} >1$ and $T_{n} > 0$. Since $B_{n+1}<0$%
, we shall have that $T_{n+1}>T_{n}>0$ provided we show that $F(-n-1,%
\mbox{$\frac{3}{2}$},\mbox{$\frac{1}{2}$};2,\mbox{$\frac{1}{4}$};1)<0$. To
this end, we again apply Identity 3 to observe that 
\[
F\left( -n-1,\mbox{$\frac{3}{2}$},\mbox{$\frac{1}{2}$};2,\mbox{$\frac{1}{4}$}%
;1\right) =\frac{\left( -\frac{1}{4}\right) _{n+1}}{\left( \frac{1}{4}%
\right) _{n+1}}F\left( -n-1,\mbox{$\frac{1}{2}$},\mbox{$\frac{1}{2}$};2,%
\mbox{$\frac{1}{4}$}-n;1\right) . 
\]
Since 
$$
\frac{\left( -\mbox{$\frac{1}{4}$}\right) _{n+1}}{\left( %
\mbox{$\frac{1}{4}$}\right) _{n+1}}<0,
$$
we need to show that 
\[
F\left( -n-1,\mbox{$\frac{1}{2}$},\mbox{$\frac{1}{2}$};2,\mbox{$\frac{1}{4}$}%
-n;1\right) >0. 
\]
Letting $m=n+1$, $a=b=\mbox{$\frac{1}{2}$}$, and $\epsilon =%
\mbox{$\frac{1}{4}$}$, it follows from Lemma \ref{required} that 
\[
F\left( -n-1,\mbox{$\frac{1}{2}$},\mbox{$\frac{1}{2}$};2,\mbox{$\frac{1}{4}$}%
-n;1\right) =F(-m,a,b;a+b+1,1+\epsilon -m;1)>0. 
\]
Hence $T_{n+1}>T_{n}>0$ for all integers $n\geq 8$ by induction.
%\thinspace \rule{2.2mm}{3.3mm}
\qquad\end{proof}

{\it Proof of Theorem~$\ref{strong}$}.
Clearly, 
$$
A_{n}=\frac{\left( \frac{1}{2}\right) _{n}\left( -\frac{1}{2}\right) _{n}}{n!n!}. 
$$
Computing the logarithmic derivative of $g$ we have 
\[
\frac{g^{\prime }(x)}{g(x)}=-\frac{1}{2}\left( \frac{(1-x)^{-\frac{1}{4}}}{%
1+(1-x)^{\frac{3}{4}}}\right), 
\]
which implies 
\begin{equation}
\left( \sum\limits_{n=0}^{\infty }(n+1)a_{n+1}x^{n}\right) \left( (1-x)^{%
\frac{1}{4}}+1-x\right) =-\frac{1}{2}\sum\limits_{n=0}^{\infty }a_{n}x^{n}.
\label{eq3}
\end{equation}
The coefficients of $x^{n}$ of the left-hand side of (\ref{eq3}) are obtained from the 
Cauchy product of the two terms.  Solving for $a_{n+1}$
yields (by extracting the $n$th and $(n-1)$st terms from the Cauchy
product) 
\begin{equation}
a_{n+1}=\frac{1}{2(n+1)}\left[ \left( \frac{5}{4}n-\frac{1}{2}\right)
a_{n}-\sum\limits_{k=0}^{n-2}(k+1)a_{k+1}\frac{\left( -\frac{1}{4}\right)
_{n-k}}{(n-k)!}\right] .  \label{eq4}
\end{equation}
We now verify (\ref{eq2}) using an inductive argument.  Clearly,  the coefficients of the terms $a_{k}$ in 
(\ref{eq4}) are nonnegative.  Computation gives: 
$a_0 = A_0 = 1,\ a_1 = A_1 = -1/4, \ a_2 = A_2 = -3/64, \ a_3 = A_3 = -5/2^{8}, \ a_4 = -11/2^{10}$ and $A_4 = -175/2^{14}$.  Suppose that the inequality in (\ref{eq2}) holds for $4 \leq k \leq n$.  From (\ref{eq4}) we have
\begin{eqnarray}
\ \ \ &&a_{n+1}\leq \label{eq5}\\
&& \frac{1}{2(n+1)}\left[ \left( \frac{5}{4}n-\frac{1}{2}\right) 
\frac{\left( \frac{1}{2}\right) _{n}\left( -\frac{1}{2}\right) _{n}}{n!n!}%
-\sum\limits_{k=0}^{n-2}(k+1)\frac{\left( \frac{1}{2}\right) _{k+1}\left( -%
\frac{1}{2}\right) _{k+1}}{(k+1)!(k+1)!}\frac{\left( -\frac{1}{4}\right)
_{n-k}}{(n-k)!}\right] .  \nonumber
\end{eqnarray}
We need to show that the right-hand side of (\ref{eq5}) is less than or equal
to $A_{n+1}=\frac{\left( \frac{1}{2}\right) _{n+1}\left( -\frac{1}{2}\right)
_{n+1}}{(n+1)!(n+1)!},$
that is,  
\begin{eqnarray}
\left( \frac{5}{4}n - \frac{1}{2}\right) \frac{\left( \frac{1}{2}\right)
_{n}\left( -\frac{1}{2}\right) _{n}}{n!n!}-2(n+1)\frac{\left( \frac{1}{2}%
\right) _{n+1}\left( -\frac{1}{2}\right) _{n+1}}{(n+1)!(n+1)!}  \label{eq6}  \\ 
\leq \sum\limits_{k=0}^{n-2}(k+1)\frac{\left( \frac{1}{2}\right) _{k+1}\left( -%
\frac{1}{2}\right) _{k+1}}{(k+1)!(k+1)!}\frac{\left( -\frac{1}{4}\right)
_{n-k}}{(n-k)!}.  \nonumber
\end{eqnarray}
After adding the $(n-1)$st and $n$th terms of the right-hand side of (%
\ref{eq6}) to inequality (\ref{eq6}) and then simplifying, we use $%
(a)_{k+1}=(a+k)(a)_{k}$, $(a)_{n-k}=(-1)^{k}(a)_{k}/(1-a-n)_{k}$, the fact that $(-n)_k = 0$ for $k \geq n+1$, and the
definition of $_{3}F_{2}$ to obtain 
\[
\frac{\left( \frac{1}{2}\right) _{n}\left( -\frac{1}{2}\right) _{n}}{n!n!}%
\cdot \frac{(2n-1)}{4(n+1)}\leq -\left( \frac{1}{4}\right) \left( -\frac{1}{4%
}\right) _{n}\frac{_{3}F_{2}\left( -n,\frac{1}{2},\frac{3}{2};2,\frac{5}{4}%
-n;1\right) }{n!},  \label{eq7}
\]
or equivalently 
\begin{equation}
_{3}F_{2}\left( -n,\mbox{$\frac{1}{2}$},\mbox{$\frac{3}{2}$};2,%
\mbox{$\frac{5}{4}$}-n;1\right) \geq \frac{\left( \frac{1}{2}\right) _{n}^{2}%
}{\left( -\frac{1}{4}\right) _{n}(n+1)!}.  \label{eq8}
\end{equation}
Clearly, the right-hand side of (\ref{eq8}) is negative for all $n\geq 1$. 
Inequality (\ref{eq8}) can be explicitly verified for $0\leq n\leq 7$. For 
$n\geq 8$, inequality (\ref{eq8}) follows from Corollary \ref{integers}.  Thus, the inequality in (\ref{eq2}) also holds for $k = n+1$.
Hence, by induction (\ref{eq2}) holds for all $k \in {\Bbb N} \cup \{0\}$.

Finally, the convexity and monotonicity of $f$ are clear.  By l'H\^{o}pital's rule, 
$f(0^{+}) = A_4 - a_4 = 1/2^{14} = 1/16384$, while the value of $f(1)$ is clear.  
%\rule{2.2mm}{3.3mm}
\qquad\endproof

\section*{Acknowledgments}

The authors would like to thank the referees for their many helpful 
suggestions which helped to considerably improve the paper.


\begin{thebibliography}{99}

\bibitem{AS}  
{\sc M. Abramowitz and I. Stegun, eds.,} 
{\em Handbook of Mathematical Functions}, 
Dover, New York, 1965.

\bibitem{AB}
{\sc G. Almkvist and B. Berndt,} 
{\it Gauss, Landen, Ramanujan, the
arithmetic-geometric mean, ellipses, pi, and the Ladies Diary,} 
Amer. Math. Monthly, 95 (1988), pp. 585--608.

\bibitem{AVV}
{\sc G. D. Anderson, M. K. Vamanamurthy, and M. Vuorinen,}
{\em Conformal Invariants, Inequalities, and Quasiconformal Mappings},
John Wiley, New York, 1997.

\bibitem{2}
{\sc R. Askey, G. Gasper, and M. Ismail,}
{\em A positive sum from summability theory,} 
J. Approx. Theory, 13 (1975), pp. 413--420.

\bibitem{BPR}
{\sc R. W. Barnard, K. Pearce, and K. C. Richards,} 
{\em A monotonicity property involving ${}_{3}F_{2}$ and comparisons 
of classical approximations of elliptical arc length,} 
SIAM J. Math. Anal., to appear.

\bibitem{BB}
{\sc J. M. Borwein and P. B. Borwein,}
{\em Pi and the AGM---A Study in Analytic Number Theory and 
Computational Complexity}, 
John Wiley, New York, 1987.

\bibitem{C}
{\sc B. C. Carlson,} 
{\em Special Functions of Applied Mathematics}, 
Academic Press, New York, 1977.

\bibitem{GR}
{\sc G. Gasper and M. Rahman,} 
{\em Basic Hypergeometric Series}, 
Cambridge University Press, Cambridge, 1990.

\bibitem{G}
{\sc A. G. Greenhill,} 
{\em The Applications of Elliptic Functions}, 
Dover, New York, 1954.

\bibitem{H}
{\sc P. Henrici,} 
{\em Applied and Computational Complex Analysis, {\rm Vol. I},} 
John Wiley, New York, 1974.

\bibitem{L}
{\sc D. F. Lawden,} 
{\em Ellipitic Functions and Applications}, 
Appl. Math. Sci. 80, Springer-Verlag, New York, 1989.

\bibitem{R}
{\sc E. Rainville,} 
{\em Special Functions}, 
Macmillan, New York, 1960.

\bibitem{V}
{\sc M. Vuorinen,} 
{\em Hypergeometric functions in geometric function theory,}
in Proceedings of the Special Functions and Differential Equations,
K.~R. Srinivasa, R. Jagannthan, and G. Van der Jeugy, eds.,
Allied Publishers, New Delhi, 1998.

\end{thebibliography}
\end{document}


